\documentclass[11pt,a4paper]{article}
\usepackage[italian]{babel}
\usepackage[margin=2.5cm]{geometry}
\usepackage{parskip}
\usepackage{amsmath, amssymb, amsthm}
\usepackage{booktabs}
\usepackage{array}
\usepackage{xcolor}
\usepackage{needspace}
\usepackage{enumitem}
\usepackage{listings}
\usepackage{tikz}
\usepackage[most]{tcolorbox}
\usepackage[hidelinks]{hyperref}
\usetikzlibrary{decorations.pathreplacing, shapes.geometric, arrows.meta, positioning}

% Niente vedove/orfane: i paragrafi non lasciano righe isolate a fine/inizio pagina
\widowpenalty=10000
\clubpenalty=10000

% Elenchi più compatti
\setlist{itemsep=2pt, parsep=0pt, topsep=4pt, partopsep=0pt}

% --- Riquadri con bordo nero, senza numero (si spezzano tra due pagine) ---
% Uso: \begin{definizione} ... \end{definizione}
%      \begin{teorema}[Nome] ... \end{teorema}  (il nome è facoltativo)
\tcbset{nota/.style={enhanced jigsaw, breakable, colback=white,
  colframe=black, boxrule=0.5pt, sharp corners}}
\NewTColorBox{definizione}{}{nota,
  before upper={\textbf{Definizione.}\ }}
\NewTColorBox{teorema}{o}{nota,
  before upper={\textbf{Teorema\IfValueT{#1}{ (#1)}.}\ }}
\NewTColorBox{proposizione}{o}{nota,
  before upper={\textbf{Proposizione\IfValueT{#1}{ (#1)}.}\ }}
\NewTColorBox{proprieta}{o}{nota,
  before upper={\textbf{Proprietà\IfValueT{#1}{ (#1)}.}\ }}
\NewTColorBox{assioma}{o}{nota, boxrule=1pt,
  before upper={\textbf{Assioma\IfValueT{#1}{ (#1)}.}\ }}

% --- Ambienti semplici (senza riquadro) ---
% Il titolo sta in cima al blocco, anche se il contenuto inizia con un riquadro o
% con codice; e non resta mai da solo in fondo a una pagina.
% Uso: \begin{esempio}[Titolo facoltativo] ... \end{esempio}
\newcommand{\newrunin}[2]{%
  \NewDocumentEnvironment{#1}{o}{%
    \par\Needspace{4\baselineskip}\addvspace{\medskipamount}\noindent
    \textbf{#2}\IfValueT{##1}{ (##1)}\textbf{.}\ \ignorespaces}%
  {\par\addvspace{\medskipamount}}}
\newrunin{esempio}{Esempio}
\newrunin{esercizio}{Esercizio}
\newrunin{osservazione}{Osservazione}

% V e F nelle tavole di verità
\newcommand{\V}{\textsf{V}}
\newcommand{\F}{\textsf{F}}

% --- Codice C ---
% Uso: \begin{codice} ... \end{codice}      codice C numerato
%      \begin{comando} ... \end{comando}    comandi da terminale
%      \begin{risultato} ... \end{risultato} output di un programma
%      \lstinline|printf("ciao");|           codice dentro al testo
\lstdefinestyle{c}{
  language=C,
  basicstyle=\ttfamily\small,
  keywordstyle=\color{blue}\bfseries,
  commentstyle=\color{gray}\itshape,
  stringstyle=\color{red!70!black},
  numbers=left,
  numberstyle=\tiny\color{gray},
  numbersep=8pt,
  columns=flexible,
  keepspaces=true,
  breaklines=true,
  showstringspaces=false,
  tabsize=4,
  morekeywords={include, define},
  % lettere accentate nei commenti
  literate={à}{{\`a}}1 {è}{{\`e}}1 {é}{{\'e}}1 {ì}{{\`i}}1
           {ò}{{\`o}}1 {ù}{{\`u}}1 {È}{{\`E}}1 {À}{{\`A}}1
}
\lstset{style=c}
\newtcblisting{codice}{listing only, nota, left=6mm,
  listing options={style=c}}
\newtcblisting{comando}{listing only, nota,
  listing options={basicstyle=\ttfamily\small, numbers=none, language={},
    columns=flexible, keepspaces=true, breaklines=true}}
\newtcblisting{risultato}{listing only, enhanced jigsaw, breakable,
  colback=black!5, colframe=black, boxrule=0.5pt, sharp corners,
  listing options={basicstyle=\ttfamily\small, numbers=none, language={},
    columns=flexible, keepspaces=true, breaklines=true}}

% --- Diagrammi di flusso (simboli standard) ---
\tikzset{
  terminale/.style = {ellipse, draw, minimum width=2.5cm, minimum height=0.9cm},
  io/.style        = {trapezium, trapezium left angle=70, trapezium right angle=110,
                      draw, minimum width=2.5cm, minimum height=0.9cm},
  processo/.style  = {rectangle, draw, minimum width=2.5cm, minimum height=0.9cm},
  decisione/.style = {diamond, draw, aspect=2, minimum width=2.5cm},
  freccia/.style   = {thick, -{Stealth}}
}

\title{Limiti notevoli di seno e coseno e conseguenze del teorema del confronto}
\author{Marco Cerbella}
\date{Lezione 14 -- 10 novembre 2026}

\begin{document}
\maketitle

\section{Limiti notevoli di seno e coseno}

\begin{teorema}[Limiti notevoli]
\[
  \lim_{x\to 0}\frac{\sin x}{x} = 1, \qquad \lim_{x\to 0}\frac{1 - \cos x}{x^2} = \frac{1}{2}.
\]
\end{teorema}

Il primo limite è stato dimostrato nella lezione precedente. Per il secondo si ha una forma indeterminata $\frac{0}{0}$:
\[
  \lim_{x\to 0}\frac{1 - \cos x}{x^2} = \frac{1 - 1}{0} = \frac{0}{0}.
\]

\begin{esempio}[Calcolo di $\lim_{x\to 0}\frac{1-\cos x}{x^2}$]
Si moltiplica numeratore e denominatore per $1 + \cos x$, usando il prodotto notevole $(a - b)(a + b) = a^2 - b^2$:
\[
  \lim_{x\to 0}\frac{1 - \cos x}{x^2}\cdot\frac{1 + \cos x}{1 + \cos x} = \lim_{x\to 0}\frac{1 - \cos^2 x}{x^2}\cdot\frac{1}{1 + \cos x}.
\]
Dall'identità fondamentale $\cos^2 x + \sin^2 x = 1$ si ricava $\sin^2 x = 1 - \cos^2 x$, quindi
\[
  = \lim_{x\to 0}\frac{\sin^2 x}{x^2}\cdot\frac{1}{1 + \cos x} = \lim_{x\to 0}\left(\frac{\sin x}{x}\right)^{2}\cdot\frac{1}{1 + \cos x}
  = \lim_{x\to 0}\frac{\sin x}{x}\cdot\frac{\sin x}{x}\cdot\frac{1}{1 + \cos x}.
\]
I tre fattori tendono rispettivamente a $1$, $1$ e $\frac{1}{2}$. Poiché il limite del prodotto è il prodotto dei limiti,
\[
  \lim_{x\to 0}\frac{1 - \cos x}{x^2} = 1\cdot 1\cdot\frac{1}{2} = \frac{1}{2}.
\]
\end{esempio}

\section{Conseguenze del teorema del confronto}

\begin{teorema}[Corollario del teorema del confronto]
Sia $f$ una funzione \emph{limitata} in un intorno di $c$, cioè esistono un intorno $I(c)$ e una costante $C > 0$ tali che
\[
  |f(x)| \le C, \quad \forall x \in I(c)\setminus\{c\}
\]
(equivalentemente $-C \le f(x) \le C$), e sia $g$ una funzione tale che $\displaystyle\lim_{x\to c} g(x) = 0$. Allora
\[
  \lim_{x\to c} f(x)\,g(x) = 0.
\]
\end{teorema}

Cioè: \emph{limitata per infinitesima è infinitesima}.

\begin{esempio}[$\lim_{x\to\pm\infty}\frac{\sin x}{x}$]
Si ha una situazione in cui il numeratore non ha limite ($\lim_{x\to+\infty}\sin x$ non esiste) e il denominatore tende a $\pm\infty$: non si può usare il rapporto dei limiti. Si scrive
\[
  \frac{\sin x}{x} = \sin x\cdot\frac{1}{x}.
\]
$\sin x$ è una funzione limitata ($-1 \le \sin x \le 1$) e $\lim_{x\to\pm\infty}\frac{1}{x} = 0$. Per il corollario del teorema del confronto,
\[
  \lim_{x\to\pm\infty}\sin x\cdot\frac{1}{x} = 0.
\]
Infatti, per $x \to +\infty$,
\[
  -\frac{1}{x} \le \frac{\sin x}{x} \le \frac{1}{x},
\]
con le due funzioni esterne che tendono a $0$ (e quindi anche quella centrale).
\end{esempio}

\begin{teorema}[Secondo teorema del confronto, caso infinito]
Siano date due funzioni $f : D_f \to \mathbb{R}$ e $g : D_g \to \mathbb{R}$. Supponiamo che
\[
  \lim_{x\to c} f(x) = +\infty
\]
e che esista un intorno $I(c)$ in cui entrambe $f$ e $g$ siano definite tranne al più il punto $c$ e tale che
\[
  f(x) \le g(x) \qquad \forall x \in I(c)\setminus\{c\}.
\]
Allora si ha anche
\[
  \lim_{x\to c} g(x) = +\infty.
\]
Un risultato analogo vale nel caso del limite $-\infty$ (con la disuguaglianza opposta).
\end{teorema}

\begin{esempio}[$\lim_{x\to+\infty}(x + \sin x)$]
$\lim_{x\to+\infty}\sin x$ non esiste, mentre $\lim_{x\to+\infty} x = +\infty$: non si può usare la somma dei limiti. Poiché $\sin x \ge -1$,
\[
  x + \sin x \ge x - 1 .
\]
Per $x \to +\infty$ si ha $x - 1 \to +\infty$; per il secondo teorema del confronto
\[
  \lim_{x\to+\infty}(x + \sin x) = +\infty.
\]
\end{esempio}

\end{document}
